\section{A common source: load demand and convex capability}\label{sec:universal}

\Cref{sec:faithful} fits a source to one prescribed lens. A machine that forms
many lenses must instead keep its cell, liquids and source layout fixed and
select the lens through the drive command alone (and the fill pumps, which are
part of the machine). This section separates the two halves of that question:
what a prescribed face \emph{demands} from the field, and what a fixed source
\emph{can} supply. The first is a closed-form property of the lens and the
liquids; the second is a convex program. Both are posed in the declared
axisymmetric, quiescent, creeping-flow reduction.

\subsection{Load demand of a prescribed face}
For a face $h(r)$ pinned at $r=R$, with tension $\sigma$ and density step
$\Delta\rho=\rho_{\text{below}}-\rho_{\text{above}}$, stationarity of the graph
energy with fixed enclosed volumes gives
\begin{equation}
 -\sigma\,\frac1r\bigl(r\,\phi\bigr)'+\Delta\rho\,g\,h=t(r)+P,\qquad
 \phi=\frac{h'}{\sqrt{1+h'^2}},
 \label{eq:demand}
\end{equation}
where $t$ is the extra normal traction (positive upward) and $P$ the face's
pressure multiplier. The \emph{demand} of a target face is the left side of
\cref{eq:demand} evaluated on the target; the field has to supply it up to an
additive constant, which the volume constraint absorbs. In the volume-preserving
surface space the demand is the mechanical gradient at the projected target,
$\bm d=\nabla_q E(q^\star)$, and an acoustic force $\bm F$ holds the target iff
$\bm F=\bm d$.

With the sag $h=h_0+\tfrac12 c\,r^2+a_4r^4+a_6r^6+\dots$, expanding $\phi$ to
fifth order gives
\begin{equation}
 t(r)=\text{const}+Q_2r^2+Q_4r^4+\dots,\qquad
 \begin{aligned}
 Q_2&=-\sigma\,(16a_4-2c^3)+\tfrac12\Delta\rho\,g\,c,\\
 Q_4&=-\sigma\,\bigl(36a_6-36c^2a_4+\tfrac94c^5\bigr)+\Delta\rho\,g\,a_4 .
 \end{aligned}
 \label{eq:demand-series}
\end{equation}
Both vanish for a sphere ($a_4=c^3/8$, $a_6=c^5/16$) with $\Delta\rho=0$, as
they must. For a conic face, $a_4=(1+k)c^3/8$, the leading term vanishes when
\begin{equation}
 \frac{\Delta\rho\,g}{\sigma}=4\,k\,c^2 ,
 \label{eq:matching-rule}
\end{equation}
so a liquid pair can match the fourth-order departure of one face from a
sphere; the sixth- and higher-order remainder is then the whole demand. The
two faces have independent outer liquids, hence independent conditions.

\paragraph{Checks.} On the published singlets of \cref{sec:faithful}'s study
S08, the generalized force of the exact traction \cref{eq:demand} and the
mechanical gradient agree to $2\times10^{-3}$; the finite-element equilibrium
under the exact traction reproduces both target faces to $3\times10^{-11}$\,m;
\cref{eq:demand-series} matches the exact traction to $6\times10^{-4}$ within
$r\le0.15R$; and \cref{eq:matching-rule} predicts the brine and oil choice that
a brute-force density and tension scan found for Fig.~9b
($-8.3\times10^4$ versus $-7.9\times10^4$\,m$^{-2}$ front, $1.0\times10^6$
versus $7.5\times10^5$\,m$^{-2}$ back).

\subsection{Convex capability of a fixed source}
About $q^\star$, a single tone with modal amplitudes $c$ exerts
$F_i=c^{\mathsf H}G_ic$ with Hermitian $G_i$ (\cref{sec:faithful}). Several tones
at distinct frequencies $\omega_t$, spaced far beyond the creeping-flow
bandwidth, exert in the time average the sum of their forces: cross terms
oscillate at $\omega_t-\omega_s$ and the interface, a low-pass filter with
relaxation rates of order $10\,\mathrm s^{-1}$, does not follow them. With
$X_t=\sum c_{t}c_{t}^{\mathsf H}$ for the tones near band $t$,
\begin{equation}
 F_i=\sum_t\operatorname{tr}\bigl(G_{t,i}X_t\bigr),\qquad X_t\succeq0,
 \qquad \text{power}=\sum_t\operatorname{tr}\bigl(W_tX_t\bigr),
 \label{eq:multitone}
\end{equation}
is linear in the positive semidefinite matrices $X_t$: the forces reachable by
multi-tone drives form a convex cone. With the frozen pupil response
$\delta\bm h=O(\bm F-\bm d)+V\bm v$, $O=HK^{-1}$, and pump trims $\bm v$ acting
through the volume columns $V$, the best pupil tolerance within a power cap,
\begin{equation}
 \min_{X_t\succeq0,\ \bm v}\ \norm{O(\bm F(X)-\bm d)+V\bm v}_\infty
 \quad\text{s.t.}\quad \sum_t\operatorname{tr}(W_tX_t)\le P_{\max},
 \label{eq:capability}
\end{equation}
is a semidefinite program. Its optimum is a global lower bound for any drive of
the given bands, single-tone or not. When it is not small, no command of that
source layout can hold the lens, and the remedy is more source freedom (bands,
modes, placement), not more iterations. When it is small, the rank of the
optimal $X_t$ says how many tones near each band the bound needs; the
single-tone case is $\operatorname{rank}X=1$, and an extreme optimal point has
rank $r$ with $r^2\le n_c$ for $n_c$ real equality constraints
\cite{barvinok1995,pataki1998}. A local rank-one refinement from the leading
eigenvector then gives a single-tone command.

\paragraph{Why this matters.} The single-tone faithful inverse of
\cref{sec:faithful} is a non-convex problem in $c$ and was started from frozen
least-squares fits of random seeds. For the published Fig.~9b singlet in the
water cell it stalled at 0.67\,\textmu m for 1\,W and for 10\,W alike, which was
read as a reachability limit of water. \Cref{eq:capability} reports instead a
pupil error below $10^{-11}$\,m at 0.46\,W; its leading eigenvector refines to a
single tone of 0.17\,W, and the full nonlinear start-up branch from the
acoustics-off meniscus under that tone ends within $5\times10^{-11}$\,m of both
faces. The earlier limit was a local minimum of the optimizer, not a property
of the liquids. The convex program is therefore used both as a certificate of
impossibility and as the initializer of the faithful inverse.

\paragraph{Scope.} The capability is a first-order, frozen-compliance statement
about the target state; it ignores the acoustic shape stiffness
($\norm{A}/\norm{K}\sim10^{-4}$ here) and the start-up path. Acceptance of a
command remains the nonlinear chain of \cref{sec:faithful}: start-up branch,
fixed-command refinement, coupled stability and formation.
